\documentclass[12pt, a4paper]{article}

\usepackage{amsmath, amssymb, amsthm}
\usepackage{geometry}
\usepackage{hyperref}
\usepackage{titlesec}
\usepackage{parskip}
\usepackage{xcolor}

\geometry{margin=2.5cm}

\hypersetup{
    colorlinks=true,
    linkcolor=blue!60!black,
    urlcolor=blue!60!black
}

% Theorem environments
\newtheorem{theorem}{Theorem}
\newtheorem{proposition}[theorem]{Proposition}
\newtheorem{lemma}[theorem]{Lemma}
\newtheorem{corollary}[theorem]{Corollary}
\theoremstyle{definition}
\newtheorem{definition}{Definition}
\newtheorem{example}{Example}
\theoremstyle{remark}
\newtheorem*{remark}{Remark}
\newtheorem*{observation}{Observation}

% Title formatting
\titleformat{\section}{\large\bfseries}{}{0em}{}[\titlerule]
\titleformat{\subsection}{\normalsize\bfseries\itshape}{}{0em}{}

\title{\textbf{On the Distinction Between Zero and Infinity}}
\author{}
\date{}

\begin{document}

\maketitle
\tableofcontents
\vspace{1cm}

\noindent\textit{A self-directed mathematical exploration examining why zero and infinity are fundamentally distinct, approached through multiple lenses: real analysis, contradiction, limit theory, and set-theoretic foundations.}

\vspace{0.5cm}
\hrule
\vspace{0.8cm}

%----------------------------------------------
\section{First Pass: What is Zero?}
%----------------------------------------------

Before proving anything, we must understand what zero actually \textit{is}. As it turns out, mathematics gives us more than one definition of zero, each revealing a different facet of its nature.

\subsection{Zero in $\mathbb{R}$: The Additive Identity}

In the real number system $\mathbb{R}$, zero is defined as the \textbf{additive identity}:
\[
\forall\, x \in \mathbb{R},\quad x + 0 = x.
\]
This property uniquely determines $0$: if there were another element $0'$ satisfying $x + 0' = x$ for all $x$, then $0 = 0 + 0' = 0'$. So zero is the \emph{unique} additive identity of $\mathbb{R}$.

\subsection{Zero in Peano Arithmetic}

In the axiomatic foundation of the natural numbers via the Peano axioms, zero is defined as the \textbf{first natural number} — the unique natural number with no predecessor. All other natural numbers are successors of some natural number; zero alone has none.

\subsection{Zero in Set Theory (Von Neumann Construction)}

In the Von Neumann construction of the natural numbers, zero is identified with the \textbf{empty set}:
\[
0 := \emptyset.
\]
Each subsequent natural number is then defined as the set of all preceding ones:
\[
1 := \{0\} = \{\emptyset\},\quad 2 := \{0,1\} = \{\emptyset, \{\emptyset\}\},\quad 3 := \{0,1,2\},\quad \ldots
\]
In this construction, zero is literally nothingness — the set containing no elements. This is perhaps the most philosophically striking definition: \emph{zero is empty sets all the way down}.

%----------------------------------------------
\section{Method 1: The Real Number System}
%----------------------------------------------

\begin{theorem}
$0 \neq \infty$ in the real number system $\mathbb{R}$.
\end{theorem}

\begin{proof}
In $\mathbb{R}$, every real number is finite and well-defined. The symbol $\infty$ is not an element of $\mathbb{R}$ — it does not belong to the set of real numbers. Since equality is defined only between elements of the same set, the expression $0 = \infty$ is not even a well-formed mathematical statement within $\mathbb{R}$: we have $0 \in \mathbb{R}$ but $\infty \notin \mathbb{R}$, so they cannot be compared as equal.

Therefore, $0 \neq \infty$.
\end{proof}

\begin{remark}
One might object: what about the \textbf{extended real line} $\overline{\mathbb{R}} = \mathbb{R} \cup \{+\infty, -\infty\}$? Here $\infty$ is formally adjoined as a symbol with a defined order: for every real $x$, we declare $x < \infty$. Since $0 \in \mathbb{R}$, we have $0 < \infty$, which immediately gives $0 \neq \infty$ in $\overline{\mathbb{R}}$ as well.
\end{remark}

%----------------------------------------------
\section{Method 2: Proof by Contradiction}
%----------------------------------------------

\begin{theorem}
$0 \neq \infty$.
\end{theorem}

\begin{proof}
Assume for contradiction that $0 = \infty$.

By the definition of $\infty$ (as used in the extended real line), for every real number $M$:
\[
\infty > M.
\]
Taking $M = 0$, this gives $\infty > 0$.

But if $0 = \infty$, then substituting:
\[
0 > 0,
\]
which is false (no real number is strictly greater than itself, by the irreflexivity of $<$).

This is a contradiction. Therefore our assumption was false, and $0 \neq \infty$.
\end{proof}

%----------------------------------------------
\section{Method 3: Limit Theory (Infinity as a Direction)}
%----------------------------------------------

The most illuminating argument may be this: infinity is not a \emph{number} — it is a \emph{direction}, a description of unbounded growth. The following makes this precise.

\begin{definition}[Limit at Infinity]
We write $\displaystyle\lim_{x \to \infty} f(x) = L$ to mean: for every $\varepsilon > 0$, there exists $M$ such that
\[
x > M \implies |f(x) - L| < \varepsilon.
\]
The notation ``$x \to \infty$'' means $x$ \textbf{becomes arbitrarily large} — it is a description of a process, not a substitution of a value.
\end{definition}

\begin{proposition}
$\displaystyle\lim_{x \to \infty} \frac{1}{x} = 0$, but this does \emph{not} mean $\dfrac{1}{\infty} = 0$.
\end{proposition}

\begin{proof}
We prove the limit rigorously. Given $\varepsilon > 0$, choose $M = \dfrac{1}{\varepsilon}$. Then for any $x > M$:
\[
\left|\frac{1}{x} - 0\right| = \frac{1}{x} < \frac{1}{M} = \varepsilon.
\]
Hence $\displaystyle\lim_{x \to \infty} \frac{1}{x} = 0$.

Crucially, in this entire proof we never substituted $x = \infty$. We only analyzed the \emph{behavior} of $\frac{1}{x}$ as $x$ grows without bound. The symbol $\infty$ here describes a \emph{direction of unbounded growth}, not a number one can substitute.
\end{proof}

\begin{observation}
This demonstrates that ``$x \to \infty$'' means $x$ becomes arbitrarily large. \textbf{Infinity is a direction of unbounded growth. It is not an element of $\mathbb{R}$.} Hence $0 \neq \infty$.
\end{observation}

%----------------------------------------------
\section{A Richer Example: Rates of Decay}
%----------------------------------------------

To see this more clearly, consider:
\[
\lim_{x \to \infty} \frac{x}{x+1}.
\]
Dividing numerator and denominator by $x$:
\[
\lim_{x \to \infty} \frac{1}{1 + \frac{1}{x}}.
\]
Since $\frac{1}{x} \to 0$, this limit equals $\dfrac{1}{1+0} = 1$.

Again, we never substituted $x = \infty$. We analyzed behavior. This distinction is fundamental.

\subsection{Which function tends to zero faster: $\frac{1}{x}$ or $\frac{1}{x^2}$?}

We say $f(x)$ goes to $0$ faster than $g(x)$ if $\displaystyle\lim_{x\to\infty} \frac{f(x)}{g(x)} = 0$.

\begin{proposition}
$\dfrac{1}{x^2}$ tends to $0$ faster than $\dfrac{1}{x}$.
\end{proposition}

\begin{proof}
Let $f(x) = \dfrac{1}{x^2}$ and $g(x) = \dfrac{1}{x}$. Then:
\[
\frac{f(x)}{g(x)} = \frac{1/x^2}{1/x} = \frac{1}{x^2} \cdot \frac{x}{1} = \frac{1}{x}.
\]
Taking the limit:
\[
\lim_{x \to \infty} \frac{f(x)}{g(x)} = \lim_{x \to \infty} \frac{1}{x} = 0.
\]
Therefore $\dfrac{1}{x^2}$ goes to zero faster than $\dfrac{1}{x}$.
\end{proof}

\begin{remark}
\textbf{Why does this work conceptually?} Because $x^2$ grows faster than $x$, a faster-growing denominator produces a faster-decaying fraction. This is the standard comparison criterion for rates of decay.
\end{remark}

\subsection{Formal $\varepsilon$-$M$ proof that $\displaystyle\lim_{x\to\infty} \frac{1}{x^2} = 0$}

\begin{proof}
Given $\varepsilon > 0$, we need $\left|\dfrac{1}{x^2} - 0\right| < \varepsilon$, i.e., $\dfrac{1}{x^2} < \varepsilon$ (since $\dfrac{1}{x^2} > 0$).

This requires $x^2 > \dfrac{1}{\varepsilon}$, i.e., $x > \dfrac{1}{\sqrt{\varepsilon}}$.

Choose $M = \dfrac{1}{\sqrt{\varepsilon}}$. Then for $x > M$:
\[
\frac{1}{x^2} < \frac{1}{M^2} = \varepsilon.
\]
Hence the limit is $0$.
\end{proof}

%----------------------------------------------
\section{Method 4: Set-Theoretic Infinity (Cardinality)}
%----------------------------------------------

Set theory reveals yet another face of infinity — one completely different from the analytic notion above.

\begin{definition}[Infinite Set — Dedekind]
A set $S$ is \textbf{infinite} if it can be put in one-to-one correspondence with a proper subset of itself.
\end{definition}

\begin{example}
The natural numbers $\mathbb{N}$ are infinite: the map $n \mapsto 2n$ is a bijection between $\mathbb{N}$ and $2\mathbb{N}$ (the even naturals), a proper subset. So $\mathbb{N} \sim 2\mathbb{N}$, and both have the same infinite size.
\end{example}

The cardinality of $\mathbb{N}$ is denoted $\aleph_0$ (aleph-null). This is a completely different concept of infinity from the $\infty$ of limits — it measures the \emph{size} of a set, not a direction of growth.

\begin{observation}
There are in fact \textbf{different sizes of infinity}. Cantor proved that $|\mathbb{R}| > |\mathbb{N}|$ — the real numbers are a ``larger'' infinity than the natural numbers. So infinity is not even a single thing; it fractures into an entire hierarchy.

Zero, by contrast, is unique: the empty set has exactly one cardinality, $0$.
\end{observation}

%----------------------------------------------
\section{Synthesis and Final Statement}
%----------------------------------------------

\begin{theorem}[Summary]
$0 \neq \infty$, and the two concepts belong to fundamentally different categories.
\end{theorem}

We have seen this from four angles:

\begin{enumerate}
    \item \textbf{Real analysis:} $\infty \notin \mathbb{R}$, so $0 = \infty$ is not even a well-formed statement in $\mathbb{R}$; and in $\overline{\mathbb{R}}$ we have $0 < \infty$ by definition.

    \item \textbf{Contradiction:} Assuming $0 = \infty$ leads to $0 > 0$, which violates the irreflexivity of the order on $\mathbb{R}$.

    \item \textbf{Limit theory:} $\infty$ is a \emph{direction} describing unbounded growth, not a number. It is a shorthand for a limiting process, never a value to substitute.

    \item \textbf{Set theory:} Infinite cardinalities like $\aleph_0$ are entirely different kinds of objects from the number $0 = |\emptyset|$, the cardinality of the empty set.
\end{enumerate}

\vspace{0.5cm}
\hrule
\vspace{0.8cm}

\noindent\textbf{The real question, left open:}

\begin{center}
\itshape Why does mathematics refuse to treat infinity as an ordinary number?
\end{center}

The short answer: because treating $\infty$ as an ordinary number breaks arithmetic. If $\infty - \infty$ and $\frac{\infty}{\infty}$ were defined as numbers, we could derive contradictions like $1 = 2$. Mathematics sacrifices the convenience of calling $\infty$ a number in order to preserve the consistency of the entire structure.

This is a question worth sitting with. The full answer leads into the theory of \emph{surreal numbers}, \emph{non-standard analysis} (where infinitesimals and infinite numbers are made rigorous), and the \emph{projective line} $\mathbb{P}^1$. Each is a different way mathematics has tried to carefully extend the number system to include infinite quantities — always with new rules, never the old ones unchanged.

\vspace{1cm}
\begin{center}
\textit{``The infinite! No other question has ever moved so profoundly the spirit of man.''} \\
--- David Hilbert
\end{center}

\end{document}
